\section{Planck constant and action line: derivation determinantal of the escala \texorpdfstring{\(10^{-34}\)}{decadica -34}}
\label{chap:escala-accion}
\index{Planck, constant}\index{action!Planck scale}

The construction of the action scale begins with an integer invariant. The
Smith normal form of the local dodecaphasic block produces a cokernel of
order \(16\); its norm on the constant section \(\alphaHMT\), followed by
a single action of \(\varphi\), places the output
\(\Sigma_H=\varphi\alphaHMT^{16}\) between \(10^{-34}\) and \(10^{-33}\). In
this manner, the characteristic decimal power of action is obtained before
joules and seconds are chosen. The angular branch subsequently supplies the two
oriented coordinates---before the return and after the return---and the
metrological chart finally expresses them as \(h\) and \(\hbar\).

\begin{exactbox}[title={Governing result: from the local cokernel to the action scale}]
The Smith form of \(H_4\), the order of its cokernel, and the inequalities determining the decade are exact. The complete result is \textsc{proved from explicit starting structure}: the local block, the invariant section, the multiplicative norm, a single action of \(\varphi\), and the decimal reader of \cref{chap:orden-determinantal-accion} are declared. The ablations identify the function of each rule and distinguish the local cokernel from the global cokernel of order \(4096\).
\end{exactbox}

\subsection{Smith Normal Form of \texorpdfstring{\(H_4\)}{H4} and Transport of \texorpdfstring{\(\Sigma_H=\varphi\alpha_{\mathrm{HMT}}^{16}\)}{Sigma H} to the Action Line}

For a of the three parallel dodecaphase orbits, the Parts~I--III obtain
\[
\operatorname{SNF}(H_4)=\operatorname{diag}(1,2,2,4),
\qquad
G_H:=\operatorname{coker}(H_4)
\cong\ZZ/2\ZZ\oplus\ZZ/2\ZZ\oplus\ZZ/4\ZZ,
\]
and therefore \(|G_H|=16\). Locality belongs to the reader's type: the
cokernel of the complete dodecaphasic transformation has order
\(16^3=4096\), but simultaneously records three direct copies and does not
replace the index of an orbit.

Let \(s_\alpha(g)=\alphaHMT\) be the constant section over \(G_H\). With the
local norm
\[
\mathcal N_{G_H}(s):=\prod_{g\in G_H}s(g)
\]
and a sola action of \(\varphi\) on the base, is recovered, with a
local label for facilitar the references of the Parte~V,
\begin{equation}
\label{eq:sigma-accion}
\Sigma_H
:=\varphi\,\mathcal N_{G_H}(s_\alpha)
=\varphi\alphaHMT^{16}.
\end{equation}
This is the same section defined in \cref{eq:sigma-h-local}, not a second
construction.

\begin{theorem}[Determinantal Order Transfer]
\label{thm:orden-accion}
Once the primitives of the local determinantal reader are fixed, the section in
\cref{eq:sigma-accion} satisfies
\[
\alphaHMT^{16}<10^{-34}
<\varphi\alphaHMT^{16}<10^{-33},
\qquad
\varepsilon_H:=\left\lfloor\log_{10}\Sigma_H\right\rfloor=-34.
\]
If \(\mathcal L_S\) is a real action line with basis
\(\mathcal U_S\), the assignment
\[
\tau_S:\RR_{>0}\longrightarrow\mathcal L_S,
\qquad x\longmapsto x\,\mathcal U_S,
\]
transports that order without modifying it. In particular, the decimal exponent
of the action coordinates defined below is \(-34\).
\end{theorem}

\begin{proof}
The exact inequalities and the value of the order are the content of \cref{thm:exponente-accion-local}; there they reduce to integer comparisons based on the constructed rational decimal of \(\alphaHMT\). The map \(\tau_S\) changes the codomain—from a positive scalar to a magnitude line with a declared basis—but preserves the scalar coordinate. Consequently, it does not alter its decimal order.
\end{proof}

The inherited evaluation is
\[
\Sigma_H
=1.046243838104756580184229208303089\ldots\times10^{-34},
\]
and the normalized number
\[
B_H:=10^{34}\Sigma_H
=1.046243838104756580184229208303089\ldots
\]
belongs to \([1,10)\). The factor \(10^{34}\) is reconstructed from the
order; it is not introduced as a threshold or prior normalization.

\subsection{Necessity of the local core}

The four exact ablations are
\[
\begin{array}{c|c}
\text{evaluated rule}&\text{decadic exponent}\\
\hline
\alphaHMT^{16}\ \text{without }\varphi&-35\\
\varphi\alphaHMT^{15}&-32\\
\varphi\alphaHMT^{17}&-37\\
\varphi\alphaHMT^{4096}&-8753.
\end{array}
\]
Thus, \(|G_H|=16\) fixes the exponent of \(\alphaHMT\) once the
local reader is chosen, and the inequalities then fix order \(-34\). These
controls exclude omission of \(\varphi\), neighboring exponents, and
substitution of the global cokernel for the local one. They do not constitute a
classification of all admissible readers or a proof of prior uniqueness
of the constant section, the unique action of \(\varphi\), or the
decimal chart.

\subsection{Posterior coupling with the action coordinates}

Write
\[
\varepsilon_H:=\left\lfloor\log_{10}\Sigma_H\right\rfloor=-34.
\]
Once the angular branch has constructed the pre-return character \(H_5\)
and the reconstructed coordinate
\[
\etaRet=\frac{\Cstar}{2}-\alphaAngle,
\]
both receive the same order in \(\mathcal L_S\):
\[
\hbar_{\HMT}^{\rm pre}
:=H_5\,10^{\varepsilon_H}\mathcal U_S,
\qquad
\hbar_{\HMT}^{\rm ret}
:=\etaRet\,10^{\varepsilon_H}\mathcal U_S.
\]
Numerically,
\[
\hbar_{\HMT}^{\rm pre}
=1.054571817646154469625821829433059\ldots
\times10^{-34}\mathcal U_S,
\]
\[
\hbar_{\HMT}^{\rm ret}
=1.054571816915039061133690584724300\ldots
\times10^{-34}\mathcal U_S.
\]
The datum \(H_5\) precedes the return and is not identified with \(\etaRet\).
Nor does it enter the proof of the exponent \(-34\).

The corresponding unreduced coordinates are
\[
h_{\HMT}^{\rm pre}
=2\pi H_5\,10^{-34}\mathcal U_S
=6.626070149999987926001110349899474\ldots
\times10^{-34}\mathcal U_S,
\]
\[
h_{\HMT}^{\rm ret}
=2\pi\etaRet\,10^{-34}\mathcal U_S
=6.626070145406254333510749847592879\ldots
\times10^{-34}\mathcal U_S.
\]
The comparison
\[
2\pi H_5-6.62607015
=-1.2073998889650100526\ldots\times10^{-14}
\]
is a subsequent external control on the mantissa; it is neither an equality with the
defining SI value nor an input to the reader.

\begin{center}
\begin{tabularx}{0.96\textwidth}{@{}L{0.22\textwidth}L{0.36\textwidth}Y@{}}
\toprule
Object & Coordenada obtenida & Mathematical and metrological function\\
\midrule
Escala determinantal &
\(\Sigma_H=1.0462438381047565\ldots\times10^{-34}\) &
Fix internally the decimal order of the action line.\\
Reduced prior action &
\(\hbar_{\rm HMT}^{\rm pre}=H_5\,10^{-34}\mathcal U_S\) &
Retain the angular coordinate prior to return.\\
Reduced posterior action &
\(\hbar_{\rm HMT}^{\rm ret}=\etaRet\,10^{-34}\mathcal U_S\) &
Publish the post-return coordinate on the same line.\\
Full cycle action &
\(h_{\rm HMT}^{\rm pre}=6.6260701499999879\ldots\times10^{-34}\mathcal U_S\) &
It transports local action through one phase turn by means of
\(h=2\pi\hbar\).\\
Carta SI &
\(h=6.62607015\times10^{-34}\ \mathrm{J\,s}\) &
external defining coordinate usada for the metrological recognition
\cite{BIPMSI}.\\
\bottomrule
\end{tabularx}
\end{center}

\subsection{\texorpdfstring{\(h\)}{h} and
\texorpdfstring{\(\hbar\)}{h bar}: action per full cycle and per radian}
\label{sec:accion-vuelta-radian-rev6}
\index{action!per lap}
\index{action!per radian}

The identities
\[
 E=h\nu=\hbar\omega,\qquad
 \omega=2\pi\nu,\qquad
 h=2\pi\hbar
\]
do not describe two quanta. They describe the same action in two charts. The
frequency \(\nu\) counts complete turns per unit time; the
angular frequency \(\omega\) counts phase in radians per unit time.
Consequently, \(h\) measures the action associated with the closure of one turn, and
\(\hbar\) measures the action per unit phase. The reduction is not merely an
abbreviation: it transports the global holonomy to the local generator of its
propagation.

The same distinction explains the coordinated use of degrees and radians. The complete
angular cycle of \(360^\circ\) makes the finite partitions visible
\[
 12\cdot30^\circ=4\cdot90^\circ
 =3\cdot120^\circ=360^\circ
\]
and the returns of \(270^\circ\) and \(360^\circ\). The radian, by contrast, is
the additive coordinate in which the exponential and the derivative are written without
inserting a conversion constant in each operation. The two charts are
related by
\[
 \mathbb R/\mathbb Z
 \xrightarrow{\;2\pi\;}
 \mathbb R/(2\pi\mathbb Z)
 \xrightarrow{\;180/\pi\;}
 \mathbb R/(360\mathbb Z).
\]
The composition preserves the point of the turn: the degree publishes its position
in the global calendar, and the radian linearizes its local motion.

Applied to the coordinates constructed above,
\[
 h_{\rm HMT}^{\rm pre}
 =2\pi\,\hbar_{\rm HMT}^{\rm pre},\qquad
 h_{\rm HMT}^{\rm ret}
 =2\pi\,\hbar_{\rm HMT}^{\rm ret}.
\]
The factor \(2\pi\) does not retrospectively correct a mantissa. It traverses one
complete turn with the local action. Thus an amplitude
\(\exp(\mathrm iS/\hbar)\), a commutation relation, and a holonomy
quantization compare an action with the local unit of phase, whereas the
global closure is expressed with \(h\).

\subsection{\(2\) causal branches and their confluence}

The angular branch and the determinantal branch both start from
\(\alphaHMT\), but neither retrospectively generates the other. The former
decomposes as
\[
\begin{gathered}
(\alphaHMT,\varphi;\,9,5,4,12,54,\text{ graded rule})
\longrightarrow H_5,\\
\alphaHMT\longrightarrow A=1000\alphaHMT\,{}^\circ,
\qquad
(\alphaHMT,\pi)\longrightarrow D_A,\\
\mathcal F_\pi\longrightarrow\rho_\pi=169,\qquad
(\rho_\pi,D_A,\alphaHMT)\longrightarrow\mathcal C_\pi(\alphaHMT),\\
(\pi,\alphaHMT,H_5,\mathcal C_\pi(\alphaHMT))
\longrightarrow y_*\longrightarrow\Cstar,\\
(A,\Cstar)\longrightarrow q_\pm
\longrightarrow\bigl(S_{90}(q_\pm),S_{120}(q_\pm)\bigr)
\longrightarrow K_{6\to8}
\longrightarrow\Gamma_{6\to8}.
\end{gathered}
\]
The posterior identification
\(\Gamma_{6\to8}\dashrightarrow\gamma_{\rm BI}\) belongs to the interfaz
fisica. The determinantal branch is
\[
(H_4,\alphaHMT,\varphi;\text{ reader rules})
\longrightarrow G_H
\longrightarrow\Sigma_H
\longrightarrow\varepsilon_H=-34.
\]
Only at the end is \(\varepsilon_H\) combined with \(H_5\) and \(\etaRet\) to
yield coordinates in \(\mathcal L_S\). Both an arrow
\(A\to H_5\) and any selection of \(C^*\) from SI or from
\(\gamma_{\rm BI}\) are excluded.

\subsection{Subsequent metrological chart and dimensional evaluation of the generated invariants}

Choosing a physical unit is an operation subsequent to the theorem. A
metrological chart is an isomorphism of lines of magnitudes
\[
\iota_{\rm SI}:\mathcal L_S\longrightarrow\mathcal L_{\mathrm{J\,s}},
\qquad
\iota_{\rm SI}(\mathcal U_S)=\mathrm{J\,s}.
\]
The unit belongs to the codomain of \(\iota_{\rm SI}\); the exponent
\(-34\) belongs to the relative arithmetic theorem. Merely juxtaposing a
unit symbol does not turn a dimensionless number into a
physical quantity.

\begin{statusbox}[title={Proof hierarchy for the Planck constant}]
The order \(10^{-34}\) follows from the local determinant and is fixed prior to
metrology. The mantissas \(H_5\) and \(\etaRet\) follow from the angular branch
and distinguish the orientations before and after the return. The identity
\(h=2\pi\hbar\) transports both coordinates between action per radian and action
per cycle. The map \(\iota_{\rm SI}\) then supplies the unit
\(\mathrm{J\,s}\) and permits the output to be compared with the defining constant
of the SI. This succession preserves in distinct places the internal derivation,
the dimensional realization, and the external comparison.
\end{statusbox}
